Trong quản trị rủi ro tài chính, hệ số tương quan tuyến tính truyền thống thường thất bại trong việc đo lường chính xác các biến động rủi ro, đặc biệt là khi thị trường rơi vào giai đoạn khủng hoảng. Bản chất của hệ số tương quan tuyến tính chỉ là đo lường mức độ liên kết tuyến tính trung bình giữa hai biến số. Tuy nhiên, trên thực tế, khi thị trường xảy ra khủng hoảng, các tài sản tài chính lại có xu hướng sụp đổ cùng nhau một cách mạnh mẽ hơn nhiều so với điều kiện bình thường; thậm chí tồn tại trường hợp hai biến có tương quan tuyến tính cao nhưng lại hoàn toàn không có sự phụ thuộc đuôi. Chính vì vậy, việc chỉ dựa vào tương quan tuyến tính hoặc áp đặt giả định độc lập có thể dẫn đến việc đánh giá thấp một cách nghiêm trọng xác suất xảy ra các tổn thất lớn. \\
Để giải quyết hạn chế này, lý thuyết Copula đã được xem như một công cụ mạnh mẽ thay thế cho hệ số tương quan tuyến tính nhờ vào khả năng tách biệt các hàm phân phối biên ra khỏi cấu trúc phụ thuộc. Tóm lại, hệ số tương quan tuyến tính thất bại vì nó giả định một mối quan hệ đồng nhất trên toàn bộ phân phối, trong khi rủi ro thực sự trong các cuộc khủng hoảng lại nằm ở sự hội tụ của các giá trị cực hạn tại đuôi phân phối, một đặc tính mà chỉ mô hình Copula mới có thể nắm bắt và phản ánh một cách hiệu quả.